% Cover letter — IEEE TCC submission. Build: pdflatex cover_letter_tcc.tex
\documentclass[10pt]{article}
\usepackage[margin=1in]{geometry}
\usepackage[T1]{fontenc}
\pagestyle{empty}
\setlength{\parskip}{0.45em}
\setlength{\parindent}{0pt}

\begin{document}

\today

Editor-in-Chief\\
IEEE Transactions on Cloud Computing

\textbf{Re: Submission of ``Pareto-Dominant Reinforcement Learning
for Cloud-Edge LLM Inference Scheduling'' (regular paper)}

Dear Editor-in-Chief,

I am pleased to submit the above manuscript for consideration in
\emph{IEEE Transactions on Cloud Computing}. The paper addresses
per-request scheduling of large-language-model (LLM) inference
across heterogeneous cloud-edge pools---to my knowledge the first
work to combine a deep-RL policy over such a pool with
serving-physics-aware state/reward design and joint SLO+energy
Pareto evaluation. Its contributions are: (i)~an open-source
simulator modeling five physical characteristics of LLM inference
(weight residency, KV-cache locality, continuous batching, the
memory-bandwidth decode floor, and the two-phase request
lifecycle); (ii)~an AIGC-aware PPO scheduler that Pareto-dominates
ten baselines---including a PerLLM-style constrained bandit and a
multi-objective NSGA-II search---on mean SLO attainment and energy
per token; (iii)~a 12-component ablation whose deliberately
reported null result shows the gain is a holistic system effect
rather than any single reward term, corroborated by a reward-weight
sensitivity study; and (iv)~a load-regime characterization of when
learned scheduling helps. Cloud-edge resource management and energy
efficiency are core topics of TCC, and the journal has recently
published closely related work on RL-based scheduling and LLM
inference resource management.

\textbf{Prior submission history (disclosure).} An earlier version
of this manuscript was submitted to \emph{Future Generation
Computer Systems} (FGCS-D-26-02940) and rejected. I wish to be
transparent about the circumstances: a LaTeX build failure in the
publisher's submission system left two of the three reviewers
without a readable manuscript (one reviewer stated the file was
unreadable; another's report discussed an entirely different
paper). The single substantive review was constructive, and this
submission addresses every one of its points---added notation and
hyperparameter tables, per-claim citations, terminology
clarifications, and a reward-weight sensitivity study---while going
substantially further: two recent learned baselines (CS-UCB
following PerLLM, and NSGA-II) were implemented and evaluated
across all configurations, and the related work was rewritten
around 19 additional 2024--2026 references. I am submitting a
self-contained LaTeX bundle that compiles cleanly in a fresh
environment, and I would be grateful if the system-generated PDF
could be checked before review; I am happy to supply a compiled PDF
immediately if any build issue arises.

This manuscript is original, is not under consideration elsewhere,
and has not been published in any form. I have no competing
interests to declare, and no funding was received for this work.
The submission follows the double-anonymized option; all
identifying information is confined to the separate title page and
this letter. The simulator and reproducibility manifests will be
released as open source upon acceptance, and an anonymized archive
can be provided to reviewers on request.

Thank you for considering this submission.

Sincerely,

Zhuolun Li\\
Faculty of Science and Engineering, University of Bristol\\
Bristol, United Kingdom\\
nu25406@bristol.ac.uk \quad ORCID: 0009-0004-1487-3862

\end{document}
